\section{Antecedentes: Contexto de la Agricultura Colombiana}

El sector agrícola colombiano representa un pilar fundamental de la economía nacional, contribuyendo significativamente al empleo, la seguridad alimentaria y los ingresos por exportaciones. Sin embargo, opera dentro de un marco socioeconómico complejo caracterizado por disparidades regionales, desafíos ambientales y desigualdades estructurales. Esta sección proporciona una visión integral del contexto agro colombiano, basándose en literatura reciente para situar las iniciativas de transformación digital dentro de trayectorias de desarrollo más amplias.

Colombia cuenta con zonas agroecológicas diversas, desde las regiones cafetaleras andinas hasta las llanuras tropicales del Llanos y las áreas costeras. Esta biodiversidad apoya una amplia gama de cultivos, incluyendo café, bananos, flores y diversas frutas y verduras, posicionando al país como un actor importante en los mercados agrícolas globales. No obstante, los pequeños productores, que constituyen la mayoría de los productores, enfrentan desafíos persistentes como acceso limitado al crédito, infraestructura inadecuada y vulnerabilidad al cambio climático.

Análisis económicos, como los de la Comisión Económica para América Latina y el Caribe (CEPAL, 2020), destacan el potencial transformador de las tecnologías de información y comunicación (TIC) en la agricultura. Estas tecnologías pueden mejorar la productividad, el acceso al mercado y las prácticas sostenibles, aunque su adopción sigue siendo desigual en las regiones. Estudios sobre "Retos de la Agricultura Digital" (Revista Colombiana de Investigaciones Agroindustriales, 2025) subrayan los déficits de infraestructura y las desigualdades socioeconómicas como barreras primarias para la digitalización en Colombia.

Factores culturales y generacionales complican aún más la adopción tecnológica. Investigaciones sobre la participación de caficultores en tecnologías digitales (2024) revelan que los enfoques basados en capacidades son esenciales para entender contextos locales y superar la resistencia. En regiones como Huila, donde predominan las prácticas agrícolas tradicionales, las herramientas digitales deben alinearse con normas culturales y realidades económicas para lograr una integración significativa.

La sostenibilidad ambiental es otra dimensión crítica del contexto agro colombiano. La variabilidad climática, la deforestación y la degradación del suelo amenazan la productividad a largo plazo, necesitando soluciones innovadoras. Trabajos de Eugênio Pacceli Reis da Fonseca et al. (2019) sobre sistemas de información verde para la gestión agroecosistema ofrecen modelos para integrar datos ambientales en la toma de decisiones agrícolas, promoviendo prácticas que equilibran productividad con custodia ecológica.

El departamento del Huila ejemplifica estas dinámicas contextuales, con sus suelos volcánicos fértiles apoyando cultivos diversos pero desafiados por problemas de conectividad y fragmentación de mercado. Abordar estos factores requiere enfoques holísticos que combinen innovación tecnológica con apoyo político y participación comunitaria. Al situar las intervenciones digitales dentro de este contexto más amplio, iniciativas como el Portal Agro Comercial del Huila pueden contribuir al desarrollo agrícola inclusivo y sostenible.

Perspectivas globales enriquecen esta comprensión, como se ve en estudios sobre aldeas inteligentes de Mohammad Raziuddin Chowdhury et al. (2023) y la integración digital en la agricultura europea (2025). Estos trabajos enfatizan el rol de la agricultura digital en lograr objetivos de desarrollo sostenible, proporcionando perspectivas comparativas para aplicaciones colombianas. De manera similar, exploraciones sobre blockchain en la agricultura sudamericana (2024) y mercados campesinos en Cundinamarca (2022) destacan oportunidades para comercio justo y cadenas cortas de suministro.

Los marcos institucionales también dan forma al contexto agro, con organizaciones como el Instituto Interamericano de Cooperación para la Agricultura (IICA, 2025) y la Organización de las Naciones Unidas para la Alimentación y la Agricultura (FAO, 2023) abogando por herramientas digitales en asistencia técnica y fortalecimiento rural. Estas entidades enfatizan la importancia de la construcción de capacidades y el alineamiento político para maximizar los beneficios de la adopción tecnológica.

En resumen, el contexto agrícola colombiano está marcado por un potencial inmenso atemperado por desafíos estructurales. La transformación digital debe abordar divisiones generacionales, restricciones económicas y vulnerabilidades ambientales para fomentar el crecimiento equitativo. La literatura revisada identifica brechas en aplicaciones digitales contextualizadas, particularmente en cerrar brechas digitales y promover modelos inclusivos, sentando las bases para intervenciones dirigidas como el portal propuesto.

Además, el legado histórico del conflicto en el campo colombiano ha dejado impactos duraderos en la confianza comunitaria y el compromiso institucional, lo que requiere enfoques que reconstruyan el capital social junto con la adopción tecnológica. El diseño centrado en la comunidad del portal aborda esto fomentando redes locales y gobernanza participativa, asegurando que las herramientas digitales refuercen en lugar de socavar las estructuras sociales tradicionales.

Las limitaciones de infraestructura, incluyendo redes eléctricas y de transporte poco confiables, agravan aún más los desafíos. Soluciones innovadoras como dispositivos alimentados por energía solar y centros de conectividad móvil son esenciales para extender el acceso digital a áreas remotas. Estas adaptaciones no solo mejoran la eficiencia operativa, sino que también contribuyen a la sostenibilidad energética en regiones fuera de la red.

Las dinámicas de género en la agricultura colombiana también merecen atención, ya que las mujeres a menudo desempeñan roles pivotales en la agricultura pero enfrentan barreras desproporcionadas al acceso a la tecnología. Un diseño inclusivo que incorpore características sensibles al género, como interfaces de usuario que acomoden niveles variables de alfabetización y restricciones de tiempo, puede empoderar a las productoras y promover la equidad de género en la agricultura digital.

Finalmente, la integración de sistemas de conocimiento indígena con tecnologías modernas ofrece un camino hacia la innovación culturalmente apropiada. Al valorar los conocimientos ecológicos locales y las prácticas tradicionales, las plataformas digitales pueden mejorar la resiliencia y la biodiversidad, alineándose con los esfuerzos globales de conservación. Esta perspectiva holística posiciona a la agricultura colombiana como líder en el avance tecnológico sostenible y culturalmente integrado.